\section{Circles, inversion, and straightedge--compass constructions}
\label{sec:circles}

This chapter works in the exact scalar field $K$ supplied by the arithmetic
interpretation.  Its elements have finite codes; equality, order, arithmetic,
and the nonnegative square root are the certified operations of that
interpretation.  In particular, no test of equality of arbitrary Cauchy names
is used.  A circle is represented by its centre $O\in K^2$ and its
\emph{squared} radius $\rho>0$:
\[
  \mathcal C(O,\rho)=\{X:\normsq{X-O}=\rho\}.
\]
Writing an equation of this form is an abbreviation for a predicate on point
codes, not an appeal to an infinite set of points.  A line is $P+tv$, where
$\normsq v>0$.  The quarter-turn is $\J(x,y)=(-y,x)$.  We repeatedly use
\[
 \dotp{\J u}{v}=\detp{u}{v},\qquad
 \detp{u}{\J u}=\normsq u,\qquad
 \normsq{u+tv}=\normsq u+2t\dotp uv+t^2\normsq v.
\]
Lengths denote nonnegative square roots.  Whenever a division occurs below,
its nonzero denominator is either displayed in the hypotheses or proved
positive in the proof.

\subsection{Intersections, chords, and tangents}

\begin{theorem}[Complete line--circle intersection test]
\label{circ:line-intersection}
For the line $P+tv$ and circle $\mathcal C(O,\rho)$, put
\[
 a=\normsq v>0,\qquad b=\dotp{P-O}{v},\qquad
 c=\normsq{P-O}-\rho,\qquad D=b^2-ac.
\]
If $D<0$, there is no intersection.  If $D=0$, there is exactly one,
$P-(b/a)v$.  If $D>0$, there are exactly two, with parameters
\[
 t_-=\frac{-b-\sqrt D}{a},\qquad t_+=\frac{-b+\sqrt D}{a}.
\]
In particular, a line meets a circle in at most two distinct points.
\end{theorem}
\begin{proof}
Expanding the circle equation gives $at^2+2bt+c=0$, and completing the
square gives the identity
\[
 a\bigl(\normsq{P+tv-O}-\rho\bigr)=(at+b)^2-D.       \tag{C1}
\]
A square cannot equal a negative element.  When $D=0$, a square is zero
only when its argument is zero, giving $t=-b/a$.  When $D>0$, let
$s=\sqrt D>0$.  Equation (C1) is equivalent to
$(at+b-s)(at+b+s)=0$.  The ordered field has no zero divisors, so one factor
is zero and the displayed two parameters are exhaustive.  Their difference
$2s/a$ is positive, so their points are distinct: multiplication by a
nonzero vector is injective on scalar parameters.

The perpendicular foot $H=P-(b/a)v$ satisfies
$\dotp{H-O}{v}=0$ and
\[
 \normsq{H-O}=\normsq{P-O}-\frac{b^2}{a},\qquad
 D=a\bigl(\rho-\normsq{H-O}\bigr).
\]
Thus the test is also the usual comparison between squared radius and
squared distance from the centre to the line.
\end{proof}

\begin{theorem}[Chord midpoint and perpendicular bisector]
\label{circ:chord-midpoint}
Let $A\ne B$ lie on $\mathcal C(O,\rho)$, put $M=(A+B)/2$ and $d=B-A$.
Then $\dotp{M-O}{d}=0$.  The perpendicular from $O$ to the chord line meets
it at $M$, and reflection in the perpendicular bisector interchanges $A$
and $B$ while fixing the circle.
\end{theorem}
\begin{proof}
Subtracting the two radius equations gives
\[
 0=\normsq{B-O}-\normsq{A-O}=2\dotp{M-O}{d}.
\]
For $H=A+td$ the condition $\dotp{H-O}{d}=0$ becomes
$(t-1/2)\normsq d=0$, whence $t=1/2$.  This also treats a diameter,
when $M=O$.  The perpendicular bisector is the line
$M+s\J d$, which is nondegenerate even in that case.

Its reflection is explicitly
\[
 R(X)=X-2\frac{\dotp{X-M}{d}}{\normsq d}\,d.
\]
For a vector $z$, expanding the square of
$z-2(\dotp zd/\normsq d)d$ gives $\normsq z$.
Consequently $R$ preserves squared distances; the preceding orthogonality
gives $R(O)=O$.  Finally $A=M-d/2$ and $B=M+d/2$ give $R(A)=B$ and
$R(B)=A$.  The radius predicate is therefore preserved.
\end{proof}

\begin{theorem}[Chord length and distance from the centre]
\label{circ:chord-distance}
For the chord above,
\[
 \normsq{B-A}=4\bigl(\rho-\normsq{M-O}\bigr).       \tag{C2}
\]
Two chords of the same circle have equal lengths exactly when their
midpoints have equal distances from the centre.  A longer chord has its
midpoint strictly nearer the centre.  Every chord has length at most the
diameter $2\sqrt\rho$, with equality exactly for a diameter.
\end{theorem}
\begin{proof}
Write $A-O=(M-O)-d/2$.  The cross term vanishes by
Theorem~\ref{circ:chord-midpoint}, so
$\rho=\normsq{M-O}+\normsq d/4$.  This is (C2).
Subtracting two instances proves the equality and strict comparison
statements.  Since $\normsq{M-O}\ge0$, it also proves
$\normsq d\le4\rho$.  Equality means $\normsq{M-O}=0$, which is equivalent
to $M=O$.  The nonnegative square root preserves these comparisons.
\end{proof}

\begin{theorem}[Thales' theorem, converse, and the inside--outside test]
\label{circ:thales}
Let $A\ne C$, $M=(A+C)/2$, and $\rho=\normsq{C-A}/4$.
For every $P$,
\[
 \dotp{A-P}{C-P}=\normsq{P-M}-\rho.                 \tag{C3}
\]
If $P\ne A,C$, the angle $APC$ is right exactly when $P$ belongs to the
circle with diameter $AC$.  If moreover $P$ is off the line $AC$, it is
obtuse inside that circle and acute outside it, with the algebraic
classification given by the sign of the dot product.
\end{theorem}
\begin{proof}
Put $h=(C-A)/2$ and $z=P-M$.  Then $A-P=-h-z$ and $C-P=h-z$.
Their dot product is $\normsq z-\normsq h$, proving (C3).
Perpendicularity is a zero dot product.  The sign of the cosine is the
sign of the dot product, since its two length denominators are positive.
For $P$ off the chord line the interior sine is positive, excluding the
zero and straight angles; hence (C3) gives the asserted acute and obtuse
cases as well.
No argument by passage to a limit is involved.
\end{proof}

\begin{theorem}[The tangent is perpendicular to the radius]
\label{circ:tangent}
Let $A\in\mathcal C(O,\rho)$ and $v\ne0$.
The line $A+tv$ has $A$ as its only intersection with the circle if and
only if $\dotp{A-O}{v}=0$.  Thus the unique tangent line at $A$ is
\[
 T_A=\{X:\dotp{X-O}{A-O}=\rho\}
     =\{A+t\J(A-O):t\in K\}.
\]
\end{theorem}
\begin{proof}
The line--circle equation, with $P=A$, is
\[
 t\bigl(t\normsq v+2\dotp{A-O}{v}\bigr)=0.
\]
If the dot product is zero, the only parameter is $t=0$.  If it is
nonzero, $t=-2\dotp{A-O}{v}/\normsq v$ is a second, nonzero parameter.
The perpendicular vectors to the nonzero vector $A-O$ are precisely the
scalar multiples of $\J(A-O)$: expand an arbitrary vector in the basis
$A-O,\J(A-O)$ and use the dot product.  This proves uniqueness and the
two equations for $T_A$.
\end{proof}

\begin{theorem}[A circle bounds its chords and has a supporting tangent]
\label{circ:convexity}
If $A,B$ lie on $\mathcal C(O,\rho)$ and $0\le t\le1$, then
\[
 \normsq{(1-t)A+tB-O}
   =\rho-t(1-t)\normsq{B-A}\le\rho.                \tag{C4}
\]
The inequality is strict if $A\ne B$ and $0<t<1$.
For every $X$ on the circle,
$\dotp{X-O}{A-O}\le\rho$, with equality exactly when $X=A$.
\end{theorem}
\begin{proof}
Expansion of the square on the left of (C4), using both radius equations,
gives (C4); the sign conditions then give the asserted inequalities.
For the tangent statement,
\[
 \normsq{X-A}=2\rho-2\dotp{X-O}{A-O}.
\]
This squared norm is nonnegative, and it is zero exactly when $X=A$.
Thus the circle lies in the indicated closed half-plane and touches its
bounding tangent at just one point.
\end{proof}

\subsection{Power of a point and two-circle geometry}

\begin{theorem}[Power of a point: secants, intersecting chords, and tangents]
\label{circ:power}
Define $\operatorname{Pow}(P)=\normsq{P-O}-\rho$.
Suppose the intersections of $P+tv$ with the circle have parameters
$t_1,t_2$, counted twice in the tangent case.  Their directed distances
$s_i=t_i\len v$ satisfy
\[
 s_1s_2=\operatorname{Pow}(P).                       \tag{C5}
\]
Consequently the products of the two directed secant segments agree for
any two secants through $P$.  If $P$ is between the endpoints of a chord,
the product of its two ordinary segment lengths is
$\rho-\normsq{P-O}$.  If $PT$ is tangent at $T$, then
\[
 \normsq{P-T}=\operatorname{Pow}(P).                 \tag{C6}
\]
\end{theorem}
\begin{proof}
Use the notation of Theorem~\ref{circ:line-intersection}.  Its two
parameters satisfy
\[
 t_1t_2=\frac{b^2-D}{a^2}=\frac{c}{a}.
\]
Multiplication by $a=\len v^2$ proves (C5), including $D=0$.
If $t_1\le0\le t_2$, the ordinary lengths are $-s_1,s_2$; their product
is $-\operatorname{Pow}(P)$.  The sign changes have therefore been made
explicit rather than suppressed in the usual secant diagram.

For a tangent at $T$, Theorem~\ref{circ:tangent} gives
$\dotp{P-T}{T-O}=0$.  Expanding
$P-O=(P-T)+(T-O)$ gives
$\normsq{P-O}=\normsq{P-T}+\rho$, which is (C6).
In particular, equal powers yield both the intersecting-chords theorem
and the tangent--secant theorem.
\end{proof}

\begin{theorem}[The two tangents from an exterior point]
\label{circ:external-tangents}
Put $p=P-O$ and $s=\normsq p$.  If $s>\rho$, the two and only two
points of tangency from $P$ are
\[
 T_\pm=O+\frac{\rho}{s}p
        \mathbin{\pm}\frac{\sqrt{\rho(s-\rho)}}{s}\,\J p.       \tag{C7}
\]
The tangent segments have equal length $\sqrt{s-\rho}$.
If $s<\rho$, there is no tangent from $P$ to a point of the circle.
\end{theorem}
\begin{proof}
For $T=O+u$ on the circle, tangency of $PT$ is equivalent to
$\dotp{p-u}{u}=0$, or $\dotp pu=\rho$.
The basis $p,\J p$ has squared norm $s>0$ in each direction, so this last
equation forces
$u=(\rho/s)p+z\J p$.  The radius equation becomes
$\rho^2/s+z^2s=\rho$, hence
$z^2=\rho(s-\rho)/s^2$.  There are exactly the two square-root choices in
(C7), and their radicand is positive.  This proves existence, distinctness,
and exhaustion.  Formula (C6) gives both equal tangent lengths.
If $s<\rho$, the equation $s=\rho+\normsq{P-T}$ required by any tangent
is impossible.
\end{proof}

\begin{theorem}[Radical axes and the radical centre]
\label{circ:radical}
For circles $\mathcal C(O_1,\rho_1)$ and
$\mathcal C(O_2,\rho_2)$ with $O_1\ne O_2$, their equal-power locus is
the line
\[
 2\dotp{X-O_1}{O_2-O_1}
       =\normsq{O_2-O_1}+\rho_1-\rho_2.              \tag{C8}
\]
It is perpendicular to the line of centres and contains every common
point of the circles.  Three circles with noncollinear centres have one
and only one point of equal power to all three.
\end{theorem}
\begin{proof}
Write $x=X-O_1$ and $d=O_2-O_1$.  Subtraction of powers gives
\[
 (\normsq x-\rho_1)-(\normsq{x-d}-\rho_2)
       =2\dotp xd-\normsq d-\rho_1+\rho_2.
\]
Thus equality is exactly (C8).  Its normal vector is the nonzero vector
$d$, and it contains the explicitly given point
$O_1+[(\normsq d+\rho_1-\rho_2)/(2\normsq d)]d$.
At a common point both powers are zero.

For three circles, set $d_i=O_i-O_1$ for $i=2,3$.
The equalities of powers $1=2$ and $1=3$ form two linear equations whose
coefficient determinant is $4\detp{d_2}{d_3}\ne0$.
Cramer's formula supplies their unique solution.  At that solution,
equality of powers $2=3$ follows by transitivity.

The hypotheses matter.  With collinear centres the radical axes may be
parallel or coincide.  With equal centres, unequal squared radii give an
empty equal-power locus, while identical circles give the whole plane.
None of these cases is silently described as a unique radical centre.
\end{proof}

\begin{theorem}[Complete intersection test for two circles]
\label{circ:circle-intersection}
For distinct centres put
\[
 d=O_2-O_1,\quad s=\normsq d>0,\quad
 h=O_1+\frac{s+\rho_1-\rho_2}{2s}d,\quad
 E=4s\rho_1-(s+\rho_1-\rho_2)^2.
\]
The circles have no common point if $E<0$, exactly the one point $h$ if
$E=0$, and exactly the two points
\[
 h\mathbin{\pm}\frac{\sqrt E}{2s}\,\J d              \tag{C9}
\]
if $E>0$.  With $r_i=\sqrt{\rho_i}$ and $\delta=\sqrt s$, intersection
is equivalently characterized by
$|r_1-r_2|\le\delta\le r_1+r_2$.
\end{theorem}
\begin{proof}
Every common point satisfies (C8), so in the basis $d,\J d$ it has the
form $h+t\J d$.  If $k=(s+\rho_1-\rho_2)/(2s)$, its first radius
equation is $k^2s+t^2s=\rho_1$, or $(2st)^2=E$.
The square-root cases give exactly (C9).  Conversely those points satisfy
the first radius equation and (C8), so equality of powers supplies the
second radius equation.  Distinctness and exhaustion follow as in the
line--circle theorem.

The identity
\[
 E=\bigl((r_1+r_2)^2-\delta^2\bigr)
       \bigl(\delta^2-(r_1-r_2)^2\bigr)
\]
follows by expanding both sides.  If $\delta$ lies between the displayed
nonnegative endpoints, both factors are nonnegative.  Below the lower
endpoint the second is negative and the first positive; above the upper
endpoint the first is negative and the second positive.  This proves the
equivalence using only ordered-field arithmetic.  When the centres
coincide, the circles are identical if their squared radii agree and
disjoint otherwise; that separate case requires no division by $s$.
\end{proof}

\begin{corollary}[Three common points determine the circle]
\label{circ:three-common}
Two circles with three distinct common points are identical.
\end{corollary}
\begin{proof}
If the centres differ, Theorem~\ref{circ:circle-intersection} permits at
most two common points.  If their centres agree, one common point already
forces equality of their squared radii.  The equality test for centre codes
supplies the finite case distinction.
\end{proof}

\begin{theorem}[Butterfly theorem with explicit intersection hypotheses]
\label{circ:butterfly}
Let $P,Q$ be distinct points of a circle and $M=(P+Q)/2$.
Let the two distinct chord lines $AB$ and $CD$ pass through $M$, with
$A,B,C,D$ on the circle.  Suppose that the lines $AD$ and $BC$ meet the
line $PQ$ in finite points $X,Y$, respectively, and are not parallel to it.
Then $M$ is the midpoint of $XY$.
\end{theorem}
\begin{proof}
Translate $M$ to zero and use the unit direction of $PQ$ as the horizontal
axis.  Thus $P=(-r,0)$, $Q=(r,0)$ with $r>0$.
The chord-midpoint theorem puts the centre on the vertical axis, say at
$(0,k)$.  Its squared radius is $r^2+k^2$, so the circle equation is
\[
 x^2+y^2-2ky=c_0,\qquad c_0=r^2>0.                 \tag{B1}
\]
Choose nonzero direction vectors $u,v$ for the two chord lines, and write
\[
 A=au,\quad B=bu,\qquad C=cv,\quad D=dv.
\]
The distinct-line assumption gives $\Delta=\detp uv\ne0$.
Substitution in (B1), followed by the quadratic root sum and product
identities, gives
\[
 a+b=\frac{2ku_y}{\normsq u},\quad ab=\frac{-c_0}{\normsq u},
 \qquad
 c+d=\frac{2kv_y}{\normsq v},\quad cd=\frac{-c_0}{\normsq v}.        \tag{B2}
\]
In particular both products are negative, so none of the four parameters
is zero.  Write the horizontal coordinates of the two intersections as
$x_X,x_Y$.  Solving the line equations $AD$ and $BC$ at vertical
coordinate zero gives
\[
 x_X=\frac{ad\Delta}{dv_y-au_y},\qquad
 x_Y=\frac{bc\Delta}{cv_y-bu_y}.                     \tag{B3}
\]
The stated nonparallel hypotheses say exactly that these denominators
are nonzero.  Adding (B3) over their common denominator gives numerator
\[
 \Delta\bigl(cd(a+b)v_y-ab(c+d)u_y\bigr).
\]
By (B2), the two terms inside the parentheses are both
$-2kc_0u_yv_y/(\normsq u\normsq v)$; their difference is zero.
Thus $x_X+x_Y=0$.  Their vertical coordinates are already zero, so
$X+Y=0=2M$.  In fact $x_X\ne0$, because its numerator and denominator
are nonzero, so $X,Y$ are distinct and this is an ordinary midpoint.

The familiar butterfly diagram places $X,Y$ on the chord segment.
That extra placement is unnecessary for the identity: the proof also
covers intersections on extensions of $PQ$.  It excludes exactly the
parallel cases in which the named finite intersection fails to exist.
\end{proof}

\subsection{Cyclic quadrilaterals and inscribed angles}

It is useful to multiply coordinate pairs as if they were complex numbers,
without introducing any new scalar objects:
\[
 (x,y)\star(u,v)=(xu-yv,xv+yu),\qquad
 \overline{(x,y)}=(x,-y).
\]
Pair multiplication is associative and commutative by expansion, and
\[
 \bar u\star v=(\dotp uv,\detp uv),\qquad
 \normsq{u\star v}=\normsq u\,\normsq v.             \tag{C10}
\]
A scalar on the right of a pair equation is identified with the pair
$(a,0)$.  Nonzero pairs are cancellable because their squared norms are
positive.  An oriented angle between nonzero vectors $u,v$ is represented
by the unit pair
\[
 z(u,v)=\frac{\bar u\star v}{\len u\len v}.
\]
The associated interior angle has pair
$(\operatorname{Re}z,|\operatorname{Im}z|)$.
This is the angle convention used elsewhere in the report.  It makes
equality and supplement statements exact finite algebra, with side
hypotheses recording the necessary branch information.

\begin{theorem}[A determinant test for concyclicity]
\label{circ:cyclic-determinant}
Let $A,B,C$ be noncollinear and let $\mathcal C(O,\rho)$ be their
circumcircle.  A fourth point $D$ lies on it exactly when
\[
 \det\begin{pmatrix}
 \normsq A&A_x&A_y&1\\
 \normsq B&B_x&B_y&1\\
 \normsq C&C_x&C_y&1\\
 \normsq D&D_x&D_y&1
 \end{pmatrix}=0.                                  \tag{C11}
\]
\end{theorem}
\begin{proof}
Replace the first column by the first column minus $2O_x$ times the
second, minus $2O_y$ times the third, plus
$(\normsq O-\rho)$ times the fourth.  Multilinearity and the vanishing of
determinants with two equal columns preserve the determinant.  Its new
entries are the four powers.  The first three are zero, so expansion in
this column gives
\[
 -\bigl(\normsq{D-O}-\rho\bigr)\detp{B-A}{C-A}.
\]
The last determinant is nonzero.  Cancellation therefore proves the
equivalence.  The circumcentre itself is obtained by the two linear
equations formed by subtracting radius equations, so the test uses no
search over circles.
\end{proof}

\begin{theorem}[Inscribed angle and central angle]
\label{circ:inscribed}
Let $A,B,P$ be pairwise distinct points of $\mathcal C(O,\rho)$.
Put $a=A-O$, $b=B-O$, $p=P-O$, $u=a-p$, $v=b-p$.
The exact directed-angle identity is
\[
 z(u,v)\star z(u,v)=\frac{\bar a\star b}{\rho}.       \tag{C12}
\]
In particular it expresses that the directed central angle is twice the
directed inscribed angle, with angles taken modulo a full turn.
In division-free coordinates it is the pair of polynomial identities
\begin{align*}
 \rho\bigl(\dotp uv^2-\detp uv^2\bigr)
       &=\normsq u\,\normsq v\,\dotp ab,\\
 2\rho\dotp uv\detp uv
       &=\normsq u\,\normsq v\,\detp ab.
\end{align*}
\end{theorem}
\begin{proof}
All three centre vectors have squared norm $\rho$.
Since $u=a-p$,
\[
 \bar a\star\bar p\star u
   =\bar p\star(\bar a\star a)
        -\bar a\star(\bar p\star p)
   =\rho(\bar p-\bar a)=-\rho\bar u.
\]
Likewise $\rho\bar v=-\bar b\star\bar p\star v$.
It follows that
\[
 (\bar a\star b)\star u\star\bar v
  =-\bar a\star\bar p\star u\star v
  =\rho\bar u\star v.
\]
Multiply this identity by $\bar u\star v$.  Commutativity gives
\[
 \rho(\bar u\star v)^2
    =(\bar a\star b)\normsq u\,\normsq v.
\]
Divide by the positive scalar
$\rho\normsq u\normsq v$ and use
$(\len u\len v)^2=\normsq u\normsq v$ to obtain (C12).
Taking its two coordinates gives the displayed polynomial identities.
Thus the theorem contains no tacit choice of an inverse trigonometric
function.  Which representative of a doubled angle is intended is
specified by the orientation of the relevant arcs or half-planes.
\end{proof}

\begin{theorem}[Angles in the same segment; opposite cyclic angles]
\label{circ:same-segment}
Let $A,B,P,Q$ lie on one circle, with $A\ne B$ and $P,Q\notin\{A,B\}$.
If $P,Q$ are on the same open side of the line $AB$, then the interior
angles $APB$ and $AQB$ are equal.  If they are on opposite sides, their
interior angle pairs are supplements: $(c,s)$ and $(-c,s)$.
Consequently opposite interior angles of a convex cyclic quadrilateral
are supplementary.
\end{theorem}
\begin{proof}
Three distinct points on a circle cannot be collinear by
Theorem~\ref{circ:line-intersection}.  Therefore the two oriented angle
pairs $z_P,z_Q$ have nonzero imaginary coordinates.  By (C12),
$z_P^2=z_Q^2$, and thus
\[
 (z_P-z_Q)\star(z_P+z_Q)=0.
\]
The sign of $\operatorname{Im}z_P$ is the sign of
\[
 \detp{A-P}{B-P}=\detp{B-A}{P-A},
\]
and similarly for $Q$.  On the same side these signs agree, so the
imaginary coordinate of $z_P+z_Q$ is nonzero.  Cancellation gives
$z_P=z_Q$.  On opposite sides the imaginary coordinate of $z_P-z_Q$ is
nonzero, so cancellation instead gives $z_P=-z_Q$.  Passing to absolute
imaginary coordinates yields equal interior pairs in the first case and
supplementary pairs in the second.

In a convex quadrilateral the two vertices opposite either diagonal lie
on opposite open sides of its line.  Apply the second conclusion to each
diagonal.  This proof records the side condition that a bare squared
angle identity would lose.
\end{proof}

\begin{theorem}[Converse angle criterion for a cyclic quadrilateral]
\label{circ:cyclic-converse}
Let $A\ne C$, and let $B,D$ lie off the line $AC$ and on opposite open
sides.  If the interior angle pairs of $ABC$ and $ADC$ are supplementary,
then $A,B,C,D$ lie on a circle.
\end{theorem}
\begin{proof}
Translate $A$ to zero, put $d=C-A$, and for $x=X-A$ define
\[
 F(x)=\normsq x-\dotp dx,\qquad T(x)=\detp dx.
\]
These are respectively the dot product and signed determinant of the
vectors $A-X,C-X$.  If $L_X=\len{A-X}\len{C-X}>0$, the interior pair at
$X$ is $(F(x)/L_X,|T(x)|/L_X)$.
Writing $b=B-A$ and $e=D-A$, supplementary pairs give
\[
 \frac{F(e)}{|T(e)|}=-\frac{F(b)}{|T(b)|}.
\]
Because $T(e),T(b)$ have opposite signs, this is exactly
$F(e)/T(e)=F(b)/T(b)$.  Let $k=F(b)/T(b)$ and set
\[
 o=\frac{d+k\J d}{2},\qquad \rho=\normsq o
       =\frac{\normsq d(1+k^2)}{4}>0.
\]
The circle with centre $A+o$ and squared radius $\rho$ has equation
\[
 \normsq x-\dotp dx-k\detp dx=0.                    \tag{C13}
\]
It contains $x=0,d,b$ by construction and $x=e$ by the ratio equality.
This gives the desired circle explicitly.  Every division was by a
nonzero determinant certified by the off-line hypotheses.
\end{proof}

\begin{theorem}[Tangent--chord angle theorem]
\label{circ:tangent-chord}
Let $A,B,P$ be three distinct points on a circle, and choose the tangent
ray at $A$ whose direction $t$ makes
$\detp{t}{B-A}$ have the same sign as
$\detp{A-P}{B-P}$.  Then the interior angle between that tangent ray and
the chord ray $AB$ equals the interior angle $APB$.
\end{theorem}
\begin{proof}
Translate the centre to zero, write $a=A-O$, $b=B-O$, and $v=b-a$.
For $t=\J a$,
\[
 \detp tv=-\dotp av=\rho-\dotp ab=\frac{\normsq v}{2}>0;
\]
the opposite tangent direction has negative determinant.  Exactly one of
the two rays therefore meets the prescribed sign condition.

Pair notation writes $\J a=i\star a$, where $i=(0,1)$.
As in the preceding inscribed-angle proof,
$\rho\bar v=-\bar a\star\bar b\star v$.
Using $\bar i=-i$ and $i^2=-1$, cancellation gives
\[
 z(t,v)^2=\frac{\bar t\star v}{t\star\bar v}
          =\frac{\bar a\star b}{\rho}.
\]
Here a quotient of nonzero pairs means multiplication by its explicit
inverse $\bar w/\normsq w$; the same identity holds for $-t$ because its
angle pair changes sign.  By (C12) this square equals the square of
$z(A-P,B-P)$.  The chosen imaginary signs agree and are nonzero, so the
cancellation argument of Theorem~\ref{circ:same-segment} proves equality
of the two oriented pairs.  Their interior angle pairs are consequently
equal.
\end{proof}

\begin{theorem}[Rational parametrization of a circle]
\label{circ:parameter}
Let $a=A-O$ be a known radius vector.  Every point of the circle except
$O-a$ is uniquely of the form
\[
 P(t)=O+\frac{1-t^2}{1+t^2}a+\frac{2t}{1+t^2}\J a,
       \qquad t\in K.                              \tag{C14}
\]
In particular, a circle with rational centre and a known rational point
has this parametrization of all its rational points, together with the
one omitted rational point.
\end{theorem}
\begin{proof}
The denominator $1+t^2$ is positive, and
$(1-t^2)^2+(2t)^2=(1+t^2)^2$.  The two vectors $a,\J a$ are orthogonal
with squared norm $\rho$, so (C14) satisfies the circle equation.
Conversely write $P-O=ua+v\J a$, where
\[
 u=\frac{\dotp{a}{P-O}}{\rho},\qquad
 v=\frac{\detp{a}{P-O}}{\rho},\qquad u^2+v^2=1.
\]
If $u=-1$, then $v^2=0$ and $P=O-a$.
Otherwise $t=v/(1+u)$ is defined.  From $v^2=1-u^2$ one obtains
$1+t^2=2/(1+u)$ and substitution recovers $u,v$ in (C14).
This inverse also proves uniqueness.  All formulas use field operations
alone and therefore preserve rationality in both directions.
\end{proof}

\subsection{Inversion and Ptolemy's theorem}

\begin{theorem}[Inversion: involution, radial order, and the distance identity]
\label{circ:inversion}
For $k>0$, define on points $P\ne O$
\[
 I(P)=O+\frac{k(P-O)}{\normsq{P-O}}.
\]
Then $I(I(P))=P$.  It preserves every ray from $O$, fixes precisely the
circle of squared radius $k$, and exchanges its inside and outside.
For $P,Q\ne O$,
\[
 \len{I(P)-I(Q)}
   =\frac{k\len{P-Q}}{\len{P-O}\len{Q-O}}.            \tag{C15}
\]
\end{theorem}
\begin{proof}
With $p=P-O$, the squared norm of the inverted vector is
$k^2/\normsq p>0$.  Applying the formula a second time gives
$k(kp/\normsq p)/(k^2/\normsq p)=p$.
The radial multiplier $k/\normsq p$ is positive, and comparing
$k^2/\normsq p$ with $k$ gives the fixed-circle and inside--outside
claims.

For $q=Q-O$, direct expansion gives
\begin{align*}
 \normsq{\normsq q\,p-\normsq p\,q}
  &=(\normsq q)^2\normsq p+(\normsq p)^2\normsq q
       -2\normsq p\normsq q\dotp pq\\
  &=\normsq p\normsq q\,\normsq{p-q}.
\end{align*}
Dividing by the positive squared denominator proves
\[
 \normsq{I(P)-I(Q)}
      =\frac{k^2\normsq{P-Q}}{\normsq p\normsq q}.
\]
Both sides of (C15) are nonnegative and have these equal squares, so they
are equal.
\end{proof}

\begin{theorem}[Inversion carries generalized circles to generalized circles]
\label{circ:inversion-circle}
Translate the inversion centre to zero.  A nonzero point $x$ satisfies
\[
 \lambda\normsq x+2\dotp ax+b=0
\]
if and only if $y=I(x)$ satisfies
\[
 b\normsq y+2k\dotp ay+\lambda k^2=0.              \tag{C16}
\]
Thus a circle through the inversion centre becomes a line not through it;
a line not through the centre becomes a circle through it, with that
centre point omitted.  A line through the centre remains the same line.
\end{theorem}
\begin{proof}
Substitute $x=ky/\normsq y$, which is valid by the involution theorem,
and multiply by $\normsq y>0$.  This gives (C16) exactly, so the
implication is reversible.

A circle with centre $c$ has $\lambda=1$, $a=-c$, and
$b=\normsq c-\rho$.  If it passes through zero, $b=0$ and (C16) becomes
$2\dotp cy=k$; this is a genuine line because $\normsq c=\rho>0$.
If $b\ne0$, completing the square in (C16) gives centre $kc/b$ and
squared radius $k^2\rho/b^2>0$.
A line has $\lambda=0$ and $a\ne0$.  For $b\ne0$, its image has centre
$-ka/b$ and squared radius $k^2\normsq a/b^2$, and passes through zero.
For $b=0$, (C16) is the original equation $\dotp ay=0$.
The origin is excluded throughout because inversion is undefined there;
it is not an additional image point.
\end{proof}

\begin{theorem}[Ptolemy's inequality]
\label{circ:ptolemy-inequality}
For any four points $A,B,C,D$ in the plane, allowing coincidences,
\[
 \len{C-A}\len{D-B}
 \le\len{B-A}\len{D-C}+\len{D-A}\len{C-B}.          \tag{C17}
\]
\end{theorem}
\begin{proof}
Translate $A$ to zero and write $b=B-A,c=C-A,d=D-A$.
Commutativity and distributivity of pair multiplication give the exact
identity
\[
 c\star(d-b)=b\star(d-c)+d\star(c-b).
\]
By (C10), $\len{u\star v}=\len u\len v$: both quantities are
nonnegative with the same square.  Apply the triangle inequality to the
displayed pair identity and replace each product length by this formula.
The result is (C17).  No nonzero denominator occurs, so the proof includes
all degenerate configurations.
\end{proof}

\begin{theorem}[Ptolemy's theorem with cyclic order made explicit]
\label{circ:ptolemy}
Let $A,B,C,D$ be distinct points of a circle in convex boundary order.
Equivalently, the four successive triple determinants
\[
 \detp{B-A}{C-A},\quad\detp{C-B}{D-B},\quad
 \detp{D-C}{A-C},\quad\detp{A-D}{B-D}
\]
are all positive, or all negative.  Then
\[
 \len{C-A}\len{D-B}
 =\len{B-A}\len{D-C}+\len{D-A}\len{C-B}.           \tag{C18}
\]
\end{theorem}
\begin{proof}
Translate $A$ to zero and invert with $k=1$.  Put
$b'=b/\normsq b$, $c'=c/\normsq c$, and $d'=d/\normsq d$.
The original circle passes through zero, so
Theorem~\ref{circ:inversion-circle} puts all three images on one line
not through zero.  In particular $b',d'$ are linearly independent:
if one were a scalar multiple of the other, their common nonzero value
under the line's linear functional would force that scalar to be $1$,
contradicting inversion's injectivity.

Write $c'=(1-t)b'+td'$.  Taking determinants with $b'$ and $d'$ gives
\[
 t=\frac{\detp{b'}{c'}}{\detp{b'}{d'}},\qquad
 1-t=\frac{\detp{c'}{d'}}{\detp{b'}{d'}}.
\]
Inversion divides each determinant by positive squared norms, so its
sign is unchanged.  The boundary-order hypotheses say that
$\detp bc$, $\detp bd$, and $\detp cd$ have the same nonzero sign.
Consequently $0<t<1$.  Thus $c'$ is between $b',d'$ and
\[
 \len{b'-d'}=\len{b'-c'}+\len{c'-d'}.
\]
Use (C15) for the three lengths:
\[
 \frac{\len{B-D}}{\len{B-A}\len{D-A}}
 =\frac{\len{B-C}}{\len{B-A}\len{C-A}}
  +\frac{\len{C-D}}{\len{C-A}\len{D-A}}.
\]
Multiplication by the positive product
$\len{B-A}\len{C-A}\len{D-A}$ proves (C18).
The explicit betweenness calculation is what makes the two products add
rather than subtract.
\end{proof}

\begin{theorem}[Converse to Ptolemy's theorem]
\label{circ:ptolemy-converse}
Suppose $A,B,C,D$ are distinct, $A,B,C$ are noncollinear, and equality
holds in (C17).  Then they are concyclic in one of the two boundary
orientations appearing in Theorem~\ref{circ:ptolemy}.
\end{theorem}
\begin{proof}
Invert about $A$ as in that proof, without initially assuming a circle.
Formula (C15) turns the equality into
\[
 \len{d'-b'}=\len{c'-b'}+\len{d'-c'}.
\]
Set $u=c'-b'$ and $v=d'-c'$; these vectors are nonzero by injectivity.
Squaring the equality gives $\dotp uv=\len u\len v>0$.
The two-dimensional Lagrange identity
$\normsq u\normsq v-\dotp uv^2=\detp uv^2$ therefore gives
$\detp uv=0$.  Hence
$v=(\dotp uv/\normsq u)u$, a positive multiple of $u$.
It follows that $c'=(1-t)b'+td'$ for some $0<t<1$.

The line through $b',c'$ does not pass through zero, since
$\detp{b'}{c'}=\detp bc/(\normsq b\normsq c)\ne0$.
Its inverse is a circle through zero by
Theorem~\ref{circ:inversion-circle}, containing $b,c,d$.
For completeness the boundary order can be read off without a picture.
Let $W=\detp{b'}{d'}\ne0$.  Then
\[
 \detp{b'}{c'}=tW,\qquad
 \detp{c'}{d'}=(1-t)W.
\]
The determinants $\detp bc,\detp cd,\detp bd$ consequently have the
same sign as $W$.  The remaining consecutive triple has the same sign
because
\[
 \detp{c-b}{d-b}
 =\frac{W\bigl((1-t)\normsq{b'}+t\normsq{d'}-\normsq{c'}\bigr)}{\normsq{b'}\normsq{c'}\normsq{d'}}
 =\frac{Wt(1-t)\normsq{d'-b'}}{\normsq{b'}\normsq{c'}\normsq{d'}}.
\]
All factors other than $W$ are positive.  This proves all four orientation
conditions.  The noncollinearity hypothesis excludes the additional
collinear equality configurations.
\end{proof}

\subsection{Finite straightedge--compass constructions}

\begin{theorem}[Copying a segment to a prescribed ray]
\label{circ:elementary-constructions}
Given a segment vector $u$ and a nonzero ray direction $v$ starting at
$P$, its unique copied endpoint is
\[
 P+\frac{\len u}{\len v}v.
\]
For $u\ne0$, it is constructed by intersecting that ray with the circle
centred at $P$ and having the length of the given segment as its radius.
\end{theorem}
\begin{proof}
The copied displacement has squared norm $\normsq u$ and nonnegative
ray parameter; uniqueness follows by solving $t^2\normsq v=\normsq u$
with $t\ge0$.  If $u=0$ the endpoint is simply $P$.
Otherwise the stated circle is nondegenerate, and its two intersections
with the whole line have parameters $\pm\len u/\len v$.
Exactly the positive parameter lies on the prescribed ray.
\end{proof}

\begin{theorem}[Midpoints, perpendiculars, and parallels are constructible]
\label{circ:perpendicular-construction}
The midpoint and perpendicular bisector of any nonzero segment can be
constructed using equal-radius circles.  From a line and any point $X$,
one can construct the perpendicular through $X$ and then the parallel
through $X$, with all intersections certified nondegenerate.
\end{theorem}
\begin{proof}
For distinct $U,V$, draw their circles of common radius $\len{V-U}$.
The equilateral construction of Theorem~\ref{tri:equilateral-construction}
gives their two intersections, so the joining line has direction
$\J(V-U)$ and passes through $(U+V)/2$.  It meets the original segment
line uniquely, since
$\detp{V-U}{\J(V-U)}=\normsq{V-U}>0$.
This gives both its perpendicular bisector and its midpoint.

For a line containing distinct $P,Q$, set
$R=\len{X-P}+\len{Q-P}>0$ and draw the circle of centre $X$ and radius
$R$.  Its distance from the line is at most $\len{X-P}<R$:
indeed the foot $H$ satisfies
$\normsq{X-P}=\normsq{X-H}+\normsq{H-P}$.
The line--circle discriminant is therefore positive, so there are two
distinct intersections $U,V$.  The point $X$ is equidistant from them,
hence it lies on their perpendicular bisector, as follows directly by
subtracting squared distances.  Constructing that bisector gives the
desired perpendicular even if $X$ originally lies on the line; it avoids
the degenerate attempt to join $X$ to itself as its own perpendicular
foot.  A second perpendicular through $X$ gives the parallel, since
$\J(\J v)=-v$.
The radius $R$ is obtained by copying the two lengths consecutively on a
ray, using the preceding theorem.
\end{proof}

\begin{theorem}[Straightedge--compass arithmetic]
\label{circ:construction-arithmetic}
Use the frame $(0,0),(1,0),(0,1)$.  Numbers represented on the horizontal
axis are closed under addition, subtraction,
multiplication, nonzero division, and nonnegative square root by explicit
straightedge--compass constructions.  Conversely, the coordinates of each
point obtained in a finite straightedge--compass construction are obtained
from the input coordinates by finitely many such scalar operations.
\end{theorem}
\begin{proof}
The preceding theorem supplies parallels and perpendiculars.  Projection
along lines parallel to the diagonal joining $(1,0)$ to $(0,1)$ transfers
$(a,0)$ to $(0,a)$ and back.  Perpendicular intersections therefore
construct $(a,b)$ from axis representatives $a,b$.

For addition, take the line through $(a,0)$ parallel to the line from
$(0,0)$ to $(b,1)$.  Its intersection with the horizontal line $y=1$ is
$(a+b,1)$; vertical projection gives $a+b$.  A half-turn about the origin
constructs $-b$, and the same operation gives subtraction.  One explicit
half-turn is the composition of the two reflections in the coordinate
axes, each obtained by laying off an equal perpendicular segment on the
other side.

For multiplication, the line through $(0,0)$ and $(a,1)$ meets $y=b$ at
$(ab,b)$, whose horizontal coordinate is $ab$.  For reciprocal, if
$a\ne0$, that same line meets $x=1$ at $(1,1/a)$, which gives $1/a$ by
coordinate transfer.  These intersections have respectively nonzero
vertical direction $1$ and nonzero horizontal direction $a$.

For square root, let $a\ge0$ and construct the circle whose diameter has
endpoints $(-1,0)$ and $(a,0)$.  Its centre is $((a-1)/2,0)$ and its
radius is $(a+1)/2>0$.  At its intersections with $x=0$ the equation is
\[
 y^2=\frac{(a+1)^2-(a-1)^2}{4}=a.
\]
Choose the nonnegative vertical coordinate and transfer it to the
horizontal axis.  For $a=0$ this is the single tangent intersection;
for $a>0$ it is the upper of two distinct intersections.  The branch is
decided on exact scalar codes.

For the converse, a line through known distinct points has coefficients
formed by subtraction and multiplication.  Two nonparallel lines are
solved by field operations and their nonzero determinant.  A circle
specified by a centre and a point has squared radius given by a sum of
two squares.  A line--circle intersection is given by
Theorem~\ref{circ:line-intersection}; a two-circle intersection is given
by Theorem~\ref{circ:circle-intersection}.  Both use field operations and
one nonnegative square root after their existence hypotheses have been
checked.  Identical circles or coincident lines do not specify a new
isolated intersection and are not valid construction steps.
Induction on the length of the finite construction list now proves the
converse, recording the displayed scalar expressions at each step.
This establishes the constructive field-operation characterization; it
does not assert that every algebraic number is constructible.
\end{proof}

\subsection{Arithmetic and formalization audit for this chapter}

Every displayed construction is a fixed composition of the coded scalar
operations.  The proofs have three kinds of steps: polynomial identities,
ordered-field deductions with explicit nonzero denominators, and selection
of the nonnegative square root of a certified nonnegative scalar.  The only
induction added by the final theorem is induction on a coded finite list of
construction steps.  The invariant records the complete finite expression
and its validated coordinate equations; it is checked by a primitive
recursive Boolean predicate.  Thus its induction is an instance of
quantifier-free induction after the usual Boolean encoding of bounded
conditions.

For a Lean development the small data types needed here are a point
$(x,y)$, a circle $(O,\rho,h_\rho)$, a line $(P,v,h_v)$, and the two angle
coordinates, all instantiated with the interpreted exact scalars.  The
intersection routines should return either a negative-discriminant
certificate, one point with a completeness proof, or two points with
distinctness and completeness proofs.  In particular, an existential
intersection theorem alone would not record the construction used here.

The first formalization targets are (C1), (C3), (C10), the determinant
identity in (C11), and the division-free form of (C12): they are polynomial
equalities.  The next layer proves positivity of the relevant norms and
discriminants, then validates the output formulas.  Ptolemy's inequality
uses the already established triangle inequality; the ordered equality
uses the inversion length identity and the explicit determinant proof of
betweenness.  These are genuine proof dependencies, rather than a call to
an ambient theorem about all real Euclidean spaces.  Connecting those
Lean lemmas to the object-level $\ipra$ checker still requires the coded
scalar interpretation described in the foundation chapter; a conventional
Lean proof over an ordered field alone is not that connection.
